\subsection*{The nonlocal overlap commutator: a necessary global source input (11 November 2026)}
\noindent
The backward Volterra estimate controls the spatial commutator of the
smooth transition, but the original quartic-density Hartree interaction
has an additional first-order \emph{nonlocal} commutator.
For $K(x)=|x|^{-2}$, $\mathcal N_H(q)=\kappa_H(K*|q|^4)|q|^2q$,
$0\le\eta\le1$ real, $f=u-v$, and $U=u-\eta f$, let
\[
 \mathcal C_\eta(u,v)
 =\mathcal N_H(U)-(1-\eta)\mathcal N_H(u)-\eta\mathcal N_H(v).
\]
Its exact first variation, with $g_{u,f}=4|u|^2\operatorname{Re}(\bar uf)$,
is
\[
 \lim_{\epsilon\to0}\epsilon^{-1}\mathcal C_\eta(u,u-\epsilon f)
 =\kappa_H|u|^2u\{\eta(K*g_{u,f})-K*(\eta g_{u,f})\}
 \quad\text{in }L^{6/5}(\mathbb R^3).
 \tag{NC-A}
\]
The local Fr\'echet derivative cancels; the full convolution commutator
need not vanish. The HLS bound $K*:L^{3/2}\to L^3$ and the two
polynomial difference inequalities give, with
$E=\|u\|_6+\|v\|_6$, $D=\|u-v\|_6$,
\[
 \|\mathcal C_\eta(u,v)\|_{6/5}
 \le455|\kappa_H|C_{\mathrm{HLS}}E^6D .
 \tag{NC-B}
\]
Thus a source-faithful bound needs global control of the selected
difference, not just a weak second derivative on the joining annulus.

Indeed choose smooth positive radial $u$, and nonzero $f\ge0$ supported
on $L/8<r<L/3$, with $f<u/2$. For the original
$\eta_L(r)=\chi(8r/L-4)$ set $v=u-f$. Then $U=u$ globally and
$f=f'=\Delta_r f=0$ throughout $[5L/8,L]$, whereas for
$3L/4<|x|<L$,
\[
 \mathcal C_{\eta_L}(u,v)(x)
 =\kappa_H u(x)^3
  \int_{\mathbb R^3}|x-y|^{-2}(u(y)^4-v(y)^4)\,dy
 \ne0,\qquad \kappa_H\ne0 .
 \tag{NC-C}
\]
Hence no bound of the full nonlocal defect solely by the local overlap
$\|\Delta_r(u-v)\|_{L^2([5L/8,L],r^2dr)}$ can hold for
arbitrary radial profiles, even with zero matching Cauchy data.
This does not rule out the true selected Hartree orbit: it isolates the
additional \emph{global density commutator} or global $L^6$ difference
that its PDE dynamics must produce with positive time weight and
parameter Lipschitz regularity. 
\paragraph{Explicit near/far decomposition.}
The genuine zero-Cauchy backward radial flux identity with
$g=\Delta_r(u-v)$ on $A_L=\{5L/8\le|x|\le L\}$ gives
\[
 |f'(r)|\le\frac{64}{25\sqrt3}L^{-1/2}\|g\|_{r,2},
 \qquad
 \|f\|_{L^\infty(A_L)}
 \le\frac{24}{25\sqrt3}L^{1/2}\|g\|_{r,2}.
\]
Since $|A_L|=129\pi L^3/128$, this provides the explicit
critical norm
\[
 \|f\|_{L^6(A_L)}
 \le C_{\rm ann}L\|g\|_{r,2},\qquad
 C_{\rm ann}=\frac{24}{25\sqrt3}
 \Bigl(\frac{129\pi}{128}\Bigr)^{1/6}.
 \tag{NC-D}
\]
Writing
$D_{\rm rem}=\|f\|_{L^6(\mathbb R^3\setminus A_L)}$,
the source-faithful estimate (NC-B) strengthens to
\[
 \boxed{
 \|\mathcal C_{\eta_L}(u,v)\|_{6/5}
 \le 455|\kappa_H|C_{\mathrm{HLS}}E^6
  (C_{\rm ann}L\|g\|_{r,2}+D_{\rm rem}).
 }\tag{NC-E}
\]
The counterexample (NC-C) satisfies $g=0$ yet
$D_{\rm rem}>0$; the global remote mismatch is therefore
indispensable unless the same Hartree PDE orbit supplies a
new cancellation. The $D_{\rm rem}$ time-weighted and parameter
Lipschitz budgets are not yet proved.

This is a paper-level calculation
only; strict original acceptance remains $11/13$.
